% Page 033 — page imprimée 29
% Contenu seul : le préambule et la pagination sont gérés par meyrueis.tex

ses clientes : « ... j’ai fait récemment breveter un procédé pour faire des chapeaux que la pluie
ne traverse pas » Très vite ces chapeaux imperméables représenteront 75\% de la production.

Le nombre d’ouvriers devait être réduit . Aucune pièce en notre possession n’y fait allusion.
Mais nous venons de retrouver aux Archives un recensement énigmatique :
En 1883 avait été recensé un foulon de feutre pour les chapeaux.
De 1884 à 1895 la chapellerie était taxée pour 3 ouvriers mais une note indique 2 ouvriers
foulonniers, 3 fabricants, 3 ouvriers finisseurs. Alors qu’en penser ?

Nous voyons seulement qu’Auguste avait en charge toute la partie commerciale et le courrier
avec les fournisseurs et les clients. Un de ses fils Albert visitait les clients comme voyageur de
commerce après avoir été représentant, en 1911, chez un chapelier de Paris ( maison Kohn et
Levy) 125 rue du Faubourg Montmartre. Son père Auguste le recommande à plusieurs
reprises à certains de ses clients pour qu’ils l’embauchent comme stagiaire ou représentant ?
Auguste commande des cloches en mérinos ou cashmere, des feutres avec des poils de lapin
qu’on ne peut teindre qu’avec des produits spéciaux surtout au moment de Pâques où les
commandes sont nombreuses.
Mais en ces années d’avant-guerre il se plaint « de la hausse constante des matières premières
et du prix de la main d’œuvre toujours plus élevé ».

La plupart des opérations se faisait à la main mais déjà apparaissent des machines : par
exemple en 1903 « j’ai perfectionné l’outillage destiné à faire la tournure aux chapeaux ;
vous reconnaîtrez la différence... » en 1908 « une passeuse et ses accessoires actionnée par
une machine à vapeur. » que l’on revendra en 1913 en Algérie pour en acheter une plus
moderne. On est en relation avec Thonnier à Lyon, (un célèbre fabricant d’aiguilles et de
machines). Il est question d’une machine Gutzner puis d’une Selecta avec leurs accessoires.
D’où avant la guerre de 1914 une certaine mécanisation de la production.

\textbf{On ne faisait quand même pas fortune :} ainsi Auguste écrit au sous préfet de Florac en
1911 pour lui demander une réduction d’impôts « ...nous avons un locataire pendant un mois
d’été et pendant ce temps nous faisons la cuisine sous un hangar et notre chambre est au
galetas sur des tréteaux... »
En 1912 c’est au Préfet qu’il s’adresse pour lui demander une réduction de sa patente : « ...les
commerces deviennent plus nombreux et plus importants sur la place centrale de Meyrueis et
il y a moins de monde à Pont Vieux... »

Auguste et ses enfants continuent toujours à s’occuper de la propriété et on vit beaucoup sur
les pommes de terre et les légumes du jardin, les châtaignes de Cabanals, les lapins ou le
cochon. Ainsi on achète 300 choux à repiquer et en même temps trois livres de truites.
La machine à vapeur est alimentée par le charbon de Lanuejols ( 20 à 30 quintaux) venant de
la mine des Gardies livré par Causse dit Carot, roulier à Lanuejols.

On achète à l’extérieur mais on vend aussi quand l’occasion se présente :

« ... Quand vous aurez du vin, envoyez moi une barrique, du bon... » 1903

« ... A la dernière foire mon père est sorti pour acheter le cochon... »

« Je vous envoie 6 doubles de châtaignes et une grillade pour vous. Elles sont très jolies,
plus que d’habitude car j’ai taillé tous les arbres de petite espèce pour les greffer. Si vous en
avez la vente je peux vous envoyer une autre sac ou deux... » 25 oct. 1902 à Blanc du Rozier.
